\documentclass[11pt]{article}
\usepackage[margin=1in]{geometry}
\usepackage[T1]{fontenc}
\usepackage[utf8]{inputenc}
\usepackage{lmodern}
\usepackage{microtype}
\usepackage{array}
\usepackage{booktabs}
\usepackage{enumitem}
\usepackage{titlesec}
\usepackage{xcolor}
\usepackage{hyperref}
\definecolor{bpsink}{HTML}{1C2530}
\definecolor{bpsaccent}{HTML}{2534B8}
\definecolor{bpsmuted}{HTML}{5B6675}
\definecolor{bpsenh}{HTML}{7A4BD0}
\definecolor{bpsline}{HTML}{CDD3DB}
\definecolor{bpssoft}{HTML}{F0F2F5}
\definecolor{bpswarnbg}{HTML}{FDF6E6}
\definecolor{bpswarnink}{HTML}{7A5A00}
\hypersetup{
  colorlinks=true,
  linkcolor=bpsaccent,
  urlcolor=bpsaccent,
  pdftitle={BPS Coding Scheme Dossier}
}
\setlist[itemize]{leftmargin=1.5em,itemsep=2pt,topsep=3pt}
\setlist[description]{style=nextline,leftmargin=0em,labelsep=0.6em,itemsep=4pt}
\titleformat{\section}{\color{bpsink}\large\bfseries}{}{0em}{}[{\color{bpsline}\titlerule}]
\titlespacing*{\section}{0pt}{1.1em}{0.5em}
\newcommand{\field}[1]{\nolinkurl{#1}}
\newcommand{\filepath}[1]{{\small\ttfamily\color{bpsmuted}\nolinkurl{#1}}}
\newcommand{\refbadge}{{\small\colorbox{bpsenh}{\textcolor{white}{\bfseries\ PROPOSED REFINEMENT \ }}\quad\itshape\color{bpsenh}Awaiting expert evaluation. Not yet applied to the pipeline.\par\smallskip}}
\newcommand{\schemeheader}[3]{%
  \begin{center}
    {\LARGE \bfseries\color{bpsink} #1}\\[0.45em]
    {\large #2}\\[0.35em]
    {\normalsize \itshape\color{bpsmuted} #3}
  \end{center}
  \vspace{0.4em}\hrule\vspace{0.9em}
}
\newcommand{\statusnote}[2]{%
  \begin{center}
  \fcolorbox{bpsline}{bpswarnbg}{%
    \parbox{0.94\linewidth}{\small\color{bpswarnink}\textbf{#1.}\ #2}%
  }
  \end{center}
  \vspace{0.6em}
}


\begin{document}
\schemeheader
  {Scheme 2}
  {Stage 2 Abstract-Level Structured Coding Scheme}
  {Corpus-wide abstract coding of BPS usage, domains, and typology}

\statusnote{DRAFT FOR EXPERT EVALUATION}{These coding schemes are a working draft circulated for expert evaluation. They have not been applied to a final review corpus. The current manuscript is a test run that exercised an earlier, coarser generation of these schemes. The workflow itself has since been validated end to end in two cross-provider test runs, in which three large language models from three different providers applied the abstract-level and the full-text scheme independently and their agreement was quantified per coded field. The full run on the review corpus is deliberately held until this evaluation is complete.}

\section*{At a Glance}
\begin{description}
\item[Workflow position] Abstract coding for every Stage 1 included record, before Stage 3 candidate selection.
\item[Operational mode] Structured LLM coding with archived JSON batches, deterministic vocabulary normalization, and a rule-based fallback when model output is incomplete or unavailable.
\item[Unit of analysis] One included review record, coded from title, abstract, publication types, and journal metadata only.
\item[Provenance basis] The actual Stage 2 output schema in stage2\_abstract\_coding.csv.
\item[Research questions] RQ1 (BPS operationalization); RQ3 (concepts and frameworks); SQ1 (conceptual problems); Feeds the Stage 3 candidate gate
\item[Tagline] Structured LLM-first coding with deterministic repair and rule-based fallback
\end{description}

\section*{Canonical Source Paths}
\begin{itemize}
\item \filepath{src/01_protocol/codebooks/stage2_codebook.md}
\item \filepath{src/06_review_stages/04_extraction/codebooks/stage2_codebook.csv}
\item \filepath{src/03_pipeline/bps_review/extraction/stage2.py}
\item \filepath{src/03_pipeline/bps_review/extraction/llm_stage2.py}
\item \filepath{src/06_review_stages/04_extraction/outputs/stage2_abstract_coding.csv}
\item \filepath{src/06_review_stages/04_extraction/outputs/stage2_llm_structured_coding.csv}
\item \filepath{src/06_review_stages/04_extraction/outputs/llm_stage2_structured_batches.jsonl}
\end{itemize}

\section*{Purpose}
This scheme standardizes abstract-level extraction for all eligible chronic pain reviews. It is the main corpus-wide coding layer used to describe review characteristics, classify the function of biopsychosocial language, detect biological, psychological, and social content, flag conceptual problems, and generate a provisional biopsychosocial typology for downstream synthesis.

It also sets the Stage 3 candidate gate: a record coded as musculoskeletal or as unspecified-but-not-excludable is carried forward for full-text work.


\section*{Carried-Through Metadata Fields}
{\small\itshape Descriptive fields retained from earlier stages for synthesis.}\par\smallskip
\begin{itemize}
\item record\_id, database, pmid, pmcid, doi
\item title, abstract (primary coding text)
\item year, journal, authors, country\_contact\_author
\item publication\_types (source publication-type metadata)
\item objective\_text (objective sentence extracted from the abstract)
\item screening\_status, screening\_reason (inherited from Stage 1)
\end{itemize}

\section*{Coded Fields and Controlled Values}
{\small\itshape The operational Stage 2 vocabulary. Every value carries an operational anchor, and adjacent values carry an explicit boundary rule, because that is what raises agreement between two coders.}\par\smallskip
\begin{description}
\item[\texttt{review\_type}] Evidence-synthesis design as stated in the abstract or publication metadata.
  \begin{itemize}[leftmargin=1.3em,itemsep=2pt,topsep=2pt]
  \item \textbf{\texttt{systematic review}} Explicit systematic methods (protocol, search, screening).
  \item \textbf{\texttt{meta-analysis}} Quantitative pooling of effect sizes.
  \item \textbf{\texttt{network meta-analysis}} Multiple-treatment comparison with indirect evidence.
  \item \textbf{\texttt{umbrella review}} Review of reviews.
  \item \textbf{\texttt{scoping or mapping review}} Breadth-oriented mapping of a field without pooled effects.
  \item \textbf{\texttt{rapid review}} Streamlined systematic methods.
  \item \textbf{\texttt{realist review}} Theory-driven mechanism review.
  \item \textbf{\texttt{integrative review}} Mixed evidence integration.
  \item \textbf{\texttt{narrative or expert review}} Non-systematic expert synthesis.
  \item \textbf{\texttt{other evidence synthesis}} A synthesis type not captured above.
  \item \textbf{\texttt{unclear}} Design cannot be determined from the abstract.
  \end{itemize}
{\small Boundary: prefer the most specific applicable label; meta-analysis outranks systematic review when pooling is explicit.}
\item[\texttt{objective\_category}] Primary stated purpose of the review.
  \begin{itemize}[leftmargin=1.3em,itemsep=2pt,topsep=2pt]
  \item \textbf{\texttt{conceptual}} Purpose is definitional, theoretical, or about operationalizing a model. {\small \textit{Choose when:} framework, concept, operationalize, model. \textit{Not this if:} A treatment-effect question, which is clinical.}
  \item \textbf{\texttt{clinical}} Purpose concerns treatment, management, rehabilitation, or care. {\small \textit{Choose when:} treatment, intervention, management, rehabilitation.}
  \item \textbf{\texttt{methodological}} Purpose concerns measurement, psychometrics, or assessment tools. {\small \textit{Choose when:} measurement, psychometric, instrument, assessment.}
  \item \textbf{\texttt{epidemiological}} Purpose concerns prevalence, incidence, risk factors, or associations. {\small \textit{Choose when:} prevalence, incidence, risk factor, association.}
  \item \textbf{\texttt{mixed}} Two or more purposes carry comparable weight.
  \item \textbf{\texttt{unclear}} Purpose cannot be determined.
  \end{itemize}
{\small Boundary: choose the dominant purpose; use mixed only when two purposes are genuinely co-primary.}
\item[\texttt{objective\_category\_source}] Whether the objective classification came from the structured LLM, a repaired LLM batch, or the deterministic fallback.
\item[\texttt{icd11\_pain\_category}] ICD-11 aligned pain category inferred from the abstract.
  \begin{itemize}[leftmargin=1.3em,itemsep=2pt,topsep=2pt]
  \item \textbf{\texttt{chronic secondary musculoskeletal pain}} Low back, neck, osteoarthritis, whiplash, and similar.
  \item \textbf{\texttt{chronic neuropathic pain}} Neuropathic, CRPS, radiculopathy.
  \item \textbf{\texttt{chronic cancer-related pain}} Cancer or oncology pain.
  \item \textbf{\texttt{chronic postsurgical or posttraumatic pain}} Post-surgical or post-traumatic persistent pain.
  \item \textbf{\texttt{chronic secondary headache or orofacial pain}} Headache, migraine, orofacial, TMD.
  \item \textbf{\texttt{chronic secondary visceral pain}} Visceral, abdominal, pelvic pain.
  \item \textbf{\texttt{chronic primary pain}} Fibromyalgia and other primary pain syndromes.
  \item \textbf{\texttt{mixed or unspecified chronic pain}} Chronic pain with no single dominant site.
  \item \textbf{\texttt{unclear}} Category cannot be inferred.
  \end{itemize}
{\small Abstract-based and allowed to remain unclear. Musculoskeletal and mixed-or-unspecified both feed the Stage 3 candidate gate.}
\item[\texttt{musculoskeletal\_flag}] Whether the review concerns musculoskeletal pain. Routes records into the musculoskeletal review.
  \begin{itemize}[leftmargin=1.3em,itemsep=2pt,topsep=2pt]
  \item \textbf{\texttt{yes}} Musculoskeletal focus is explicit.
  \item \textbf{\texttt{no}} Clearly non-musculoskeletal.
  \item \textbf{\texttt{unclear}} Cannot be determined; carried forward.
  \end{itemize}
\item[\texttt{neuropathic\_flag}] Whether the review concerns neuropathic pain. Parallel to musculoskeletal\_flag; routes records into the neuropathic review. A record may set both flags (for example a mixed review) or neither.
  \begin{itemize}[leftmargin=1.3em,itemsep=2pt,topsep=2pt]
  \item \textbf{\texttt{yes}} Neuropathic focus is explicit (for example neuropathy, radiculopathy, CRPS, postherpetic or diabetic neuropathic pain).
  \item \textbf{\texttt{no}} Clearly non-neuropathic.
  \item \textbf{\texttt{unclear}} Cannot be determined; carried forward.
  \end{itemize}
\item[\texttt{stage3\_track}] Which review or reviews the record is routed to. Derived from the two flags; a record can belong to both pools.
  \begin{itemize}[leftmargin=1.3em,itemsep=2pt,topsep=2pt]
  \item \textbf{\texttt{musculoskeletal}} Musculoskeletal review only.
  \item \textbf{\texttt{neuropathic}} Neuropathic review only.
  \item \textbf{\texttt{both}} Enters both pools (mixed or unresolved pain family).
  \item \textbf{\texttt{none}} Neither track; not a Stage 3 candidate.
  \end{itemize}
\item[\texttt{bps\_mention\_location}] Where the BPS term appears.
  \begin{itemize}[leftmargin=1.3em,itemsep=2pt,topsep=2pt]
  \item \textbf{\texttt{title only}} Only in the title.
  \item \textbf{\texttt{abstract only}} Only in the abstract.
  \item \textbf{\texttt{title and abstract}} In both.
  \item \textbf{\texttt{unclear}} Location cannot be resolved.
  \end{itemize}
\item[\texttt{bps\_function}] The rhetorical and analytic work the BPS label performs. This is the central RQ1 field and the finest-grained one.
  \begin{itemize}[leftmargin=1.3em,itemsep=2pt,topsep=2pt]
  \item \textbf{\texttt{explanatory framework}} BPS is used as a model that explains pain or pain-related disability. {\small \textit{Choose when:} explains, accounts for, mechanism of. \textit{Not this if:} Only justifies multimodal treatment, which is intervention rationale. \textit{Boundary:} Requires explanatory intent, not only endorsement of a model's existence.}
  \item \textbf{\texttt{intervention rationale}} BPS mainly justifies multimodal or interdisciplinary treatment. {\small \textit{Choose when:} supports multidisciplinary care. \textit{Not this if:} Explains pain aetiology, which is explanatory framework.}
  \item \textbf{\texttt{organizing principle}} BPS structures the scope or categories of the review without specifying integration mechanisms. {\small \textit{Choose when:} We organize factors into bio, psycho, social. \textit{Not this if:} A stated cross-domain mechanism.}
  \item \textbf{\texttt{justification}} BPS justifies the relevance or importance of the review topic.
  \item \textbf{\texttt{background framing}} BPS sets context in the background without analytic use.
  \item \textbf{\texttt{conclusion}} BPS appears mainly in the concluding claims.
  \item \textbf{\texttt{policy/practice implication}} BPS frames a policy or practice recommendation.
  \item \textbf{\texttt{rhetorical label}} BPS is invoked ceremonially, aspirationally, or symbolically with no substantive analytic work. {\small \textit{Choose when:} a biopsychosocial approach is needed, with no follow-through. \textit{Not this if:} Any substantive explanatory or organizing use.}
  \item \textbf{\texttt{unclear}} Function cannot be determined.
  \end{itemize}
\item[\texttt{bio\_mentioned / psych\_mentioned / social\_mentioned}] Binary presence of each domain in the abstract.
  \begin{itemize}[leftmargin=1.3em,itemsep=2pt,topsep=2pt]
  \item \textbf{\texttt{yes}} The domain is present beyond the mere BPS token.
  \item \textbf{\texttt{no}} The domain is not present.
  \end{itemize}
{\small The field is binary on purpose: a domain is either substantively present or it is not, and the mere BPS token never counts as presence.}
\item[\texttt{quality\_assessment\_reported}] Whether the abstract reports a risk-of-bias or quality assessment.
  \begin{itemize}[leftmargin=1.3em,itemsep=2pt,topsep=2pt]
  \item \textbf{\texttt{yes}} AMSTAR, risk of bias, or quality assessment named.
  \item \textbf{\texttt{no}} Not reported.
  \item \textbf{\texttt{unclear}} Cannot be determined.
  \end{itemize}
\item[\texttt{psychological\_concepts\_detected}] Normalized list of psychological concepts from title and abstract. Feeds Scheme 5. \textit{(free text)}
\item[\texttt{theoretical\_frameworks\_detected}] Normalized list of named frameworks or model labels. \textit{(free text)}
\item[\texttt{conceptual\_problem\_flags}] Provisional abstract-level conceptual problems.
  \begin{itemize}[leftmargin=1.3em,itemsep=2pt,topsep=2pt]
  \item \textbf{\texttt{vague\_definition}} BPS or a key construct is used without a definition.
  \item \textbf{\texttt{tokenistic\_bps}} BPS appears as a label with no analytic follow-through.
  \item \textbf{\texttt{missing\_social}} Social domain absent despite a BPS claim.
  \item \textbf{\texttt{missing\_biology}} Biological domain absent despite a BPS claim.
  \item \textbf{\texttt{mechanistic\_absence}} No cross-domain mechanism is offered.
  \item \textbf{\texttt{construct\_overlap}} Constructs are used interchangeably or with blurred boundaries.
  \item \textbf{\texttt{parallel\_listing\_without\_integration}} Domains are listed side by side with no integration.
  \item \textbf{\texttt{none}} No inferable conceptual problem.
  \end{itemize}
\item[\texttt{provisional\_typology}] Abstract-level BPS signal, refined later at Stage 3.
  \begin{itemize}[leftmargin=1.3em,itemsep=2pt,topsep=2pt]
  \item \textbf{\texttt{potential integrative signal}} All three domains substantively present with a cross-domain signal.
  \item \textbf{\texttt{multifactorial signal}} All three domains present but mainly parallel.
  \item \textbf{\texttt{pseudo-bps or partial signal}} One or more core domains thin or absent despite BPS language.
  \item \textbf{\texttt{rhetorical label signal}} BPS mainly symbolic or confined to framing or conclusion.
  \end{itemize}
\item[\texttt{stage3\_candidate / stage3\_priority}] Whether and how urgently the record proceeds to Stage 3.
  \begin{itemize}[leftmargin=1.3em,itemsep=2pt,topsep=2pt]
  \item \textbf{\texttt{yes / no}} Candidate flag.
  \item \textbf{\texttt{high / medium / low}} Retrieval and coding priority.
  \end{itemize}
\item[\texttt{coding\_rationale}] One-sentence rationale supporting the coding bundle. \textit{(free text)}
\item[\texttt{coding\_method / llm\_model}] Operational provenance of the row (llm\_structured, llm\_batch\_fallback, rule\_based\_fallback) and the model identifier when the LLM stage was used.
\end{description}

\section*{Prompt Rules Used in the Structured LLM Stage}
{\small\itshape These constraints are embedded in the structured prompt and are the operational meaning of the codes above.}\par\smallskip
\begin{itemize}
\item Code only from title, abstract, publication types, and journal metadata. Do not infer content that is absent from the record.
\item Do not treat the single lexical token biopsychosocial as proof that all three domains are substantively present.
\item Use explanatory framework only when BPS explicitly explains pain or pain-related disability.
\item Use intervention rationale when BPS mainly justifies multimodal or interdisciplinary treatment.
\item Use organizing principle when BPS structures scope or categories without integration mechanisms.
\item Use rhetorical label when BPS is ceremonial, aspirational, or symbolic with no substantive analytic work.
\item Set stage3\_candidate to yes for musculoskeletal or neuropathic reviews, and for mixed or unspecified chronic pain reviews where either relevance cannot be ruled out. The two planned reviews (musculoskeletal and neuropathic) draw from this same candidate pool, routed by pain-condition family.
\item Use potential integrative signal only when all three core domains are substantively present and the abstract signals cross-domain explanation or organization beyond simple listing.
\item Return conceptual\_problem\_flags as none only when no inferable conceptual problem is present.
\end{itemize}

\section*{Deterministic Repair and Fallback Logic}
{\small\itshape How the pipeline guarantees a complete, in-vocabulary row.}\par\smallskip
\begin{itemize}
\item Missing or malformed model responses are repaired against the fixed vocabularies above.
\item Aliases such as scoping review or narrative review are normalized to the implemented labels.
\item stage3\_candidate is forced to yes whenever musculoskeletal\_flag or neuropathic\_flag is yes or unclear, and stage3\_track is set from whichever flags fire (musculoskeletal, neuropathic, or both).
\item Conceptual problem flags are post-derived from typology, BPS function, and missing domains: rhetorical usage triggers tokenistic\_bps and vague\_definition; parallel non-mechanistic usage triggers parallel\_listing\_without\_integration and mechanistic\_absence; absent biological or social content triggers missing\_biology or missing\_social.
\item When the LLM stage fails, the rule-based fallback still generates review type, objective category, ICD-11 category, BPS function, domain mention flags, concept detections, and provisional Stage 3 priority.
\end{itemize}

\section*{Primary Outputs Using This Scheme}
\begin{itemize}
\item \filepath{src/06_review_stages/04_extraction/outputs/stage2_abstract_coding.csv}
\item \filepath{src/06_review_stages/04_extraction/forms/stage2_double_code_subset.csv}
\item \filepath{src/06_review_stages/04_extraction/outputs/stage2_objective_llm_assist.csv}
\item \filepath{src/06_review_stages/04_extraction/outputs/llm_objective_pilot.json}
\end{itemize}

\end{document}
